% GENERATED FILE. Edit canonical lessons, then run npm run generate:books.
% canonical-source-hash: fnv1a64:9f2910f561538621
% canonical-lessons: GE-C378-hausarbeit, GE-C378-referat, GE-C378-aufsatz, GE-C378-zusammenfassung, GE-C378-zusammenfassen

\chapter{Die Hausarbeit}
\label{ch:ge-die-hausarbeit}
\hlchaptermodality{378}

\begin{chapteropening}
\textbf{By the end of this chapter:} \emph{I can talk about written work for a course: a term paper, a talk in class, an essay and a summary.}

Complete the last lesson of chapter 378: I can talk about written work for a course: a term paper, a talk in class, an essay and a summary.
\end{chapteropening}

\section[die Hausarbeit]{die Hausarbeit — the term paper; housework}

\label{lesson:GE-C378-hausarbeit}



\begin{quote}
Before the new one: say the German for the degree programme, then the German for the internship.
\end{quote}

\subsection*{You'll want to know: die Hausarbeit}
\textbf{die Hausarbeit} — \textquotedblleft{}the term paper; housework\textquotedblright{}.

\begin{quote}
Bis Ende März muss ich eine Hausarbeit über die Industrie schreiben.
\end{quote}

\begin{quote}
\emph{By the end of March I have to write a term paper about industry.}
\end{quote}

Plural: \emph{die Hausarbeiten}. It also means housework: \emph{die Hausarbeit teilen}.

\begin{cousinweb}
\emph{Hausarbeit} joins \emph{Haus} \textquotedblleft{}house\textquotedblright{} to \emph{Arbeit} \textquotedblleft{}work\textquotedblright{}, Old High German \emph{arbeit}, Proto-Germanic \emph{arbaidiz} (EtymDB): work done at home.
\end{cousinweb}

\begin{grammarlens}[title={What you've built}]
One more word for telling, explaining and describing at B1.
\end{grammarlens}

\subsection*{Your turn}
\begin{itemize}
\raggedright
  \item \emph{Say it:} \emph{die Hausarbeit}
  \item \emph{Say it:} \emph{die Hausarbeit}, once more
  \item \emph{Say it:} a sentence of your own with \emph{die Hausarbeit}
  \item \emph{Recall:} say the German for the degree programme, then the German for the internship, then say \emph{die Hausarbeit} again
\end{itemize}

\begin{tcolorbox}[breakable,colback=teal!4,colframe=teal!35!black,title={Before you move on}]
What does die Hausarbeit mean? (\textquotedblleft{}The term paper; housework\textquotedblright{}.) Say it once more.
\end{tcolorbox}
\section[das Referat]{das Referat — the class talk, the oral report}

\label{lesson:GE-C378-referat}



\begin{quote}
Before the new one: say the German for the internship, then the German for the term paper.
\end{quote}

\subsection*{You'll want to know: das Referat}
\textbf{das Referat} — \textquotedblleft{}the class talk, the oral report\textquotedblright{}.

\begin{quote}
Morgen halte ich ein Referat über Wale.
\end{quote}

\begin{quote}
\emph{Tomorrow I am giving a talk in class about whales.}
\end{quote}

Plural: \emph{die Referate}. \emph{Ein Referat halten} — to give a talk. In an office, \emph{das Referat} is also a section.

\begin{cousinweb}
\emph{Referat} comes from the Latin verb form \emph{referat} \textquotedblleft{}let him report\textquotedblright{}, from \emph{referre} \textquotedblleft{}to carry back, to bring back word\textquotedblright{} (Lewis \& Short). It is usually said to have been first a note on a document, then the report itself.
\end{cousinweb}

\begin{grammarlens}[title={What you've built}]
One more word for telling, explaining and describing at B1.
\end{grammarlens}

\subsection*{Your turn}
\begin{itemize}
\raggedright
  \item \emph{Say it:} \emph{das Referat}
  \item \emph{Say it:} \emph{das Referat}, once more
  \item \emph{Say it:} a sentence of your own with \emph{das Referat}
  \item \emph{Recall:} say the German for the internship, then the German for the term paper, then say \emph{das Referat} again
\end{itemize}

\begin{tcolorbox}[breakable,colback=teal!4,colframe=teal!35!black,title={Before you move on}]
What does das Referat mean? (\textquotedblleft{}The class talk, the oral report\textquotedblright{}.) Say it once more.
\end{tcolorbox}
\section[der Aufsatz]{der Aufsatz — the essay, the composition}

\label{lesson:GE-C378-aufsatz}



\begin{quote}
Before the new one: say the German for the term paper, then the German for the class talk.
\end{quote}

\subsection*{You'll want to know: der Aufsatz}
\textbf{der Aufsatz} — \textquotedblleft{}the essay, the composition\textquotedblright{}.

\begin{quote}
Im Deutschunterricht haben wir einen Aufsatz über unsere Ferien geschrieben.
\end{quote}

\begin{quote}
\emph{In German class we wrote an essay about our holidays.}
\end{quote}

Plural: \emph{die Aufsätze}. Also a top part set on something: \emph{der Aufsatz} of a cupboard.

\begin{cousinweb}
\emph{Aufsatz} is usually traced to Middle High German \emph{ūfsaz}, of the family of \emph{aufsetzen} \textquotedblleft{}to set up, to draft\textquotedblright{}, built on \emph{setzen} \textquotedblleft{}to set\textquotedblright{}: a text you set down.
\end{cousinweb}

\begin{grammarlens}[title={What you've built}]
One more word for telling, explaining and describing at B1.
\end{grammarlens}

\subsection*{Your turn}
\begin{itemize}
\raggedright
  \item \emph{Say it:} \emph{der Aufsatz}
  \item \emph{Say it:} \emph{der Aufsatz}, once more
  \item \emph{Say it:} a sentence of your own with \emph{der Aufsatz}
  \item \emph{Recall:} say the German for the term paper, then the German for the class talk, then say \emph{der Aufsatz} again
\end{itemize}

\begin{tcolorbox}[breakable,colback=teal!4,colframe=teal!35!black,title={Before you move on}]
What does der Aufsatz mean? (\textquotedblleft{}The essay, the composition\textquotedblright{}.) Say it once more.
\end{tcolorbox}
\section[die Zusammenfassung]{die Zusammenfassung — the summary}

\label{lesson:GE-C378-zusammenfassung}



\begin{quote}
Before the new one: say the German for the class talk, then the German for the essay.
\end{quote}

\subsection*{You'll want to know: die Zusammenfassung}
\textbf{die Zusammenfassung} — \textquotedblleft{}the summary\textquotedblright{}.

\begin{quote}
Am Ende des Kapitels steht eine kurze Zusammenfassung.
\end{quote}

\begin{quote}
\emph{At the end of the chapter there is a short summary.}
\end{quote}

Plural: \emph{die Zusammenfassungen}.

\begin{cousinweb}
\emph{Zusammenfassung} is the noun of \emph{zusammenfassen}, built from \emph{zusammen} \textquotedblleft{}together\textquotedblright{} and \emph{fassen} \textquotedblleft{}to grasp\textquotedblright{}, Old High German \emph{fazzōn}, Proto-Germanic \emph{fatōną} (EtymDB): you grasp it all together.
\end{cousinweb}

\begin{grammarlens}[title={What you've built}]
One more word for telling, explaining and describing at B1.
\end{grammarlens}

\subsection*{Your turn}
\begin{itemize}
\raggedright
  \item \emph{Say it:} \emph{die Zusammenfassung}
  \item \emph{Say it:} \emph{die Zusammenfassung}, once more
  \item \emph{Say it:} a sentence of your own with \emph{die Zusammenfassung}
  \item \emph{Recall:} say the German for the class talk, then the German for the essay, then say \emph{die Zusammenfassung} again
\end{itemize}

\begin{tcolorbox}[breakable,colback=teal!4,colframe=teal!35!black,title={Before you move on}]
What does die Zusammenfassung mean? (\textquotedblleft{}The summary\textquotedblright{}.) Say it once more.
\end{tcolorbox}
\section[zusammenfassen]{zusammenfassen — to sum up, to summarize}

\label{lesson:GE-C378-zusammenfassen}



\begin{quote}
Before the new one: say the German for the essay, then the German for the summary.
\end{quote}

\subsection*{You'll want to know: zusammenfassen}
\textbf{zusammenfassen} — \textquotedblleft{}to sum up, to summarize\textquotedblright{}.

\begin{quote}
Die Lehrerin fasst die wichtigsten Punkte noch einmal zusammen.
\end{quote}

\begin{quote}
\emph{The teacher sums up the most important points once more.}
\end{quote}

Separable: \emph{sie fasst zusammen}, \emph{sie hat zusammengefasst}.

\begin{cousinweb}
\emph{zusammenfassen} is built from \emph{zusammen} \textquotedblleft{}together\textquotedblright{} and \emph{fassen} \textquotedblleft{}to grasp, to hold\textquotedblright{}, inherited from Old High German \emph{fazzōn} and Proto-Germanic \emph{fatōną} (EtymDB).
\end{cousinweb}

\begin{grammarlens}[title={What you've built}]
One more word for telling, explaining and describing at B1.
\end{grammarlens}

\subsection*{Your turn}
\begin{itemize}
\raggedright
  \item \emph{Say it:} \emph{zusammenfassen}
  \item \emph{Say it:} \emph{zusammenfassen}, once more
  \item \emph{Say it:} a sentence of your own with \emph{zusammenfassen}
  \item \emph{Recall:} say the German for the essay, then the German for the summary, then say \emph{zusammenfassen} again
\end{itemize}

\begin{tcolorbox}[breakable,colback=teal!4,colframe=teal!35!black,title={Before you move on}]
What does zusammenfassen mean? (\textquotedblleft{}To sum up, to summarize\textquotedblright{}.) Say it once more.
\end{tcolorbox}
